\documentclass[10pt]{beamer}
\usetheme{Antibes}
\usecolortheme{dolphin}
\usepackage[utf8]{inputenc}
\usepackage[T1]{fontenc}
\usepackage{amsmath, amssymb, amsfonts}

\title{Spectral Analysis of the Neural Tangent Kernel}
\subtitle{Review of: \textit{Eigenvalue Distribution of the NTK in the Quadratic Scaling} (Benigni \& Paquette, 2025)}
\author{Janis AIAD}
\institute{UMD}
\date{\today}

\begin{document}

% --- SLIDE 1: Title ---
\begin{frame}
    \titlepage
\end{frame}

% --- SLIDE 2: The Object of Study: NTK Dynamics ---
\begin{frame}
    \frametitle{The Object of Study: NTK Dynamics}
    \begin{block}{Gradient Flow Dynamics}
        \begin{equation*}
            \frac{d\hat{y}_t}{dt} = - \nabla_\theta \hat{y}(X, \theta_t)^\top \nabla_\theta \hat{y}(X, \theta_t) (\hat{y}_t - y) = -K^{\text{NTK}}(\theta_t) (\hat{y}_t - y)
        \end{equation*}
    \end{block}
    \vspace{1cm}
    \begin{alertblock}{Key Insight}
    The NTK spectrum governs the eigenspeeds of learning, defining the functional modes acquired during training.
    \end{alertblock}
\end{frame}

% --- SLIDE 3: Model Formulation ---
\begin{frame}
    \frametitle{Model Formulation}
    \begin{block}{2-Layer NTK at Initialization}
        \begin{equation*}
            K^{\text{NTK}} = \frac{1}{d} XX^\top + \frac{1}{p} (XX^\top) \circ \left( \phi\left(\frac{XW}{\sqrt{d}}\right) D^2 \phi\left(\frac{XW}{\sqrt{d}}\right)^\top \right)
        \end{equation*}
    \end{block}
    \begin{exampleblock}{Quadratic Scaling Regime (Assumption 2.1)}
        \begin{equation*}
            \frac{n}{d} \to \gamma_1 > 0 \quad \text{and} \quad \frac{d}{p} \to \gamma_2 > 0
        \end{equation*}
    \end{exampleblock}
    \begin{alertblock}{Key Insight}
    This scaling regime moves beyond the lazy-training limit, ensuring the non-linear term contributes non-trivially to the spectrum.
    \end{alertblock}
\end{frame}

% --- SLIDE 4: Main Result: The Limiting Spectral Distribution ---
\begin{frame}
    \frametitle{Main Result: The Limiting Spectral Distribution}
    \begin{theorem}[Thm. 2.7]
        The Empirical Spectral Distribution $L_{K^{\text{NTK}}}$ converges weakly to $\mu$, characterized by:
        \begin{equation*}
            \mu = \mu_{\text{mp}}^{\gamma_1} \boxtimes \mu_{\nu, \phi}
        \end{equation*}
    \end{theorem}
    \begin{block}{Stieltjes Transform of the Marchenko-Pastur Map (Def. 2.5)}
        For a measure $\nu$ and its Stieltjes transform $s_\nu(z)$, the transform $s(z)$ of $\mu_{\text{mp}}^\gamma \boxtimes \nu$ solves:
        \begin{equation*}
            s(z) = s_\nu\left(z - \gamma \int \frac{t}{1+ts(z)} d\mu_{\text{mp}}(t) \right)
        \end{equation*}
    \end{block}
    \begin{alertblock}{Key Insight}
    The NTK spectrum is a free multiplicative convolution of the data's Gram matrix spectrum and a novel law induced by the non-linear features.
    \end{alertblock}
\end{frame}

% --- SLIDE 5: Proof Strategy I: Decoupling via Chaos Expansion ---
\begin{frame}
    \frametitle{Proof Strategy I: Decoupling via Chaos Expansion}
    \begin{block}{Wiener-Hermite Chaos Expansion of $\phi$}
        \begin{equation*}
            \phi(x) = \mathbb{E}_{z \sim \mathcal{N}}[\phi(z)] + \alpha_\phi x + \psi(x) \quad \text{where } \alpha_\phi := \mathbb{E}[z\phi(z)]
        \end{equation*}
    \end{block}
    \begin{exampleblock}{Key Decoupling Step (Prop. 3.1)}
        Let $\tilde{K}$ be the kernel where $\psi(XW/\sqrt{d})$ is replaced by an independent noise field $\tilde{\psi}(\tilde{X})$. Then for $z \in \mathbb{C}^+$,
        \begin{equation*}
            \left| \frac{1}{n}\text{Tr}(K-zI)^{-1} - \frac{1}{n}\text{Tr}(\tilde{K}-zI)^{-1} \right| \xrightarrow[n\to\infty]{P} 0
        \end{equation*}
    \end{exampleblock}
    \begin{alertblock}{Key Insight}
    The global spectrum is sensitive only to the linear projection of the non-linearity; higher-order terms can be replaced by independent noise.
    \end{alertblock}
\end{frame}

% --- SLIDE 6: Proof Strategy II: RMT on the Covariance Tensor ---
\begin{frame}
    \frametitle{Proof Strategy II: RMT on the Covariance Tensor}
    \begin{block}{Analysis via Bai-Zhou Resolvent Method}
    The problem reduces to finding the spectrum of a conditional covariance tensor $Q$, given $W$ and $D$.
    \end{block}
    \begin{exampleblock}{Structure of the Limiting Covariance Tensor (Sec. 5)}
        \begin{equation*}
            Q = \frac{\alpha_\phi^2}{d} \left( \tau_{23}(\text{Id}_d \otimes W^\top W D) + \tau_{24}(WD \otimes WD) \right) + \beta_\phi^2 \tau_{23}(\text{Id}_d \otimes D^2)
        \end{equation*}
    \end{exampleblock}
    \begin{alertblock}{Key Insight}
    The problem is mapped from the NTK's spectrum to the spectrum of a deterministic operator constructed from the random weight matrices.
    \end{alertblock}
\end{frame}

% --- SLIDE 7: Proof Strategy III: Spectrum of the Operator Q ---
\begin{frame}
    \frametitle{Proof Strategy III: Spectrum of the Operator Q}
    \begin{block}{Spectrum of Q (Lemma 5.4)}
        The asymptotic eigenvalue distribution of the operator $Q$ is given by $\text{Law}(\frac{\alpha_\phi^2}{d}\chi + \beta_\phi^2)$, where the law of $\chi$ is:
        \begin{equation*}
            \text{Law}(\chi) = \frac{\gamma_2}{2}(\mu_{\text{mp}}^{\gamma_2} * \mu_{\text{mp}}^{\gamma_2}) + (1-\gamma_2)\mu_{\text{mp}}^{\gamma_2} + \frac{\gamma_2}{2}\delta_0
        \end{equation*}
    \end{block}
    \begin{exampleblock}{Origin of MP Law (Thm. 5.2)}
    The squared singular values of $W \in \mathbb{R}^{d \times p}$ converge to $\mu_{\text{mp}}^{\gamma_2}$.
    \end{exampleblock}
    \begin{alertblock}{Key Insight}
    The foundational Marchenko-Pastur law from the weight matrix $W$ propagates through the analysis, forming the building block of the final spectrum.
    \end{alertblock}
\end{frame}

% --- SLIDE 8: Discussion: Open Problems & Future Directions ---
\begin{frame}
    \frametitle{Discussion: Open Problems & Future Directions}
    \begin{itemize}
        \item \textbf{The Spectral Edge (Mini-Bulk):} The analysis covers the bulk density. What is the universality class of the largest eigenvalues which dictate initial learning?
        \vspace{0.5cm}
        \item \textbf{Finite-Width Corrections:} This is the leading-order term in a $1/d$ expansion. Computing the tensor-valued finite-width corrections is a major open problem.
        \vspace{0.5cm}
        \item \textbf{Scaling with Depth:} How does this spectral structure evolve in deep networks? This requires connecting with distinct RMT formalisms like those of Terjek et al.
        \vspace{0.5cm}
        \item \textbf{Structured Data:} The concentration assumption on $X$ precludes data on low-dimensional manifolds. How does data geometry impact the NTK spectrum?
    \end{itemize}
    \begin{alertblock}{Key Insight}
    This work provides a foundational result upon which more complex, practical, and finite models can be analyzed.
    \end{alertblock}
\end{frame}

% --- SLIDE 9: Conclusion ---
\begin{frame}
    \frametitle{Conclusion}
    \begin{center}
        \vspace{1.5cm}
        \Huge Thank You
    \end{center}
\end{frame}

\end{document}
